\subsection*{Jeff's rough notes}

\begin{itemize}
    \item Some positives:
    \begin{itemize}
        \item The idea of a general regression framework for non-Euclidean data in metric spaces is useful as it is an important problem in modern science.
        \item While the manifold learning aspect appears in numerous other methods in the literature, none have combined together the manfold learning with subsequent deep learning on the embeddings and with local Fr\'echet regression.
        \item The theoretical results provide a solid grounding for the framework and is one of the strengths of the work
        \item the taxi network example is interesting and a great illustration of the potential of the method
    \end{itemize}
    \item Some weaknesses/limitations/concerns:
    \begin{itemize}
        \item The focus of the method is on providing point estimates of conditional means.  This has some value, but there is no indication of how to perform statistical inference on these quantities to assess statistical significance, and given the highly flexible manifold learning and deep learning on the regressors, such inference would have to account for these sources of uncertainty so is far from trivial.
        \item The manifold learning is a critical element of the paper given that all downstream steps condition on the chosen lower dimensional embeddings space, which loses any variability not captured by those embeddings. However, the criteria for which approach to use, how many dimensions to choose, or how well those dimensions capture the salient structure of the data objects are not clarified.  
        \item While their framework could utilize any manifold learning approach the authors mention several, the one they use ISOMAP (which they claim is the only one to give good results for their examples), I believe the literature suggests that particular approach can be somewhat fragile and not robust to certain properties of the manifold.
        \item The distributional data setting used as the basis for the simulation is a practically important application of this method, but the simulation chosen is quite simplistic -- with each subject's distribution being Gaussian, so a simple mean and log variance transformed model could be easily fit using standard regression methods, and not require the more advanced tools developed here.
        \item It is not clear whether the authors share software to fit their methods, or how general the software is.
        \item The deep learning on the covariates is one of the novel steps in this approach, but it is not clear under what settings it yields or worse results than other simpler approaches, and the authors do not provide guidelines for determining when it is appropriate.  For example, are there some settings in which it could fail or add instability, and a simpler approach would be preferable? 
        \item The paper says the method accommodates high dimensional predictors, but the examples are moderate dimensional at best, on the order of a dozen or so.  While nonlinear deep-learning can detect important interactions on dimensions like these, it is not clear how large a covariate space will work well with this method -- for example if hundreds or thousands of predictors, as in genomics.  The mention of "high dimensional" may imply it applies to these settings, but it is not clear whether it did.  It would be helpful for the authors to provide a sense for how large a covariate space could be handled, and guidance on how to decide if the deep learning is appropriate for their given setting.
        \item Their examples and simulation all are low dimensional (r=2) and not very curved, so it is not clear how well the approach would perform in practice for structured objects inherently higher dimensional ($>>2$) and with more stark non-Euclidean features.
        \item For the network regression example, in fact they compute their predictions and distances in the graph Laplacian space, not the network space, so is not strictly network regression.  Multiple networks can map to the same graph Laplacian so it is not invertible.
        

    \end{itemize}
    \item Page 1438, first column paragraph 2: states "One key innovation of the proposed approach is that we demonstrate that local Fr\'echet regression can be used to map the low-dimensional manifold representation to the original metric space so that the final output resides in the metric space where the random objects are situated", and "This map, which has not yet been considered in the manifold learning literature for metric space-valued data, is crucial for the implementation of regression model in statistics, as it is mandatory for interpretability to predict the random objects residing in metric spaces, rather than their lower dimensional representations".  
    
    Is it really true that this strategy has not been considered in previous literature?

    This is the underlying idea motivating much of our previous work in the functional mixed model framework, using various transformations to a latent space for modeling, but with the inverse transformations mapping back to the original data space for interpretation and and inference, if desired.  Much of the work involved Euclidean spaces, but as described in Morris (2017 Statistical Modeling), is applicable to objects on non-Euclidean spaces using appropriate transformations/inverse transforms including manifold learning approaches that ensure the latent embeddings are mapped back to the appropriate space, and several recent works have been applied to objects on non-Euclidean manifolds (see examples below).  
    
    This work may not have explicitly used the local Fr\'echet regression modeling strategy underlying this paper, but it seems like we should point out that this is the precise motivation of this whole vein of work, which has explicitly been applied to structured functions on non-Euclidean manifolds.
    \begin{itemize}
        \item The Lee et al. (2019 JASA) paper dealt with a spherical domain, using a spherical wavelet representation to map into a set of lower-dimensional embeddings for modeling, but with the inverse spherical wavelet transform mapping back onto the sphere.
        \item The Desai et al. (2026 Biometrics) paper deals with correlation matrix objects, using the matrix logarithm (which they call "multivariate Fisher's transformation) to map to an unconstrained space for the modeling, with the inverse transform ensuring valid correlation space realizations.
        \item Morris (2017 Statistical Modeling) overviewed the general functional mixed modeling framework whose development spanned many papers starting in Morris et al. (2003 JASA), highlighting its application to functional data, potentially residing on non-Euclidean manifolds, using a suitable manifold learning to get lower dimensional embeddings, and using inverse transforms to map the embeddings back to the appropriate non-Euclidean object space.  This framework proposes a strict near-lossless criteria to guide the manifold learning to ensure accurate representation of the objects.
        \item Zohner et al. (2025 arXIV, JCGS under revision) present a framework and graphical interface to investigate the near-losslessness of various potentially invertible transformations to latent spaces, including both linear and nonlinear methods, and including manifold learning, to choose which representation seems most suitable for a given application, and to choose the latent embedding dimension to achieve the desired level of near-losslessness.
        \item We also should cite the Gunning et al. DIGITS paper (which will be on arxiv by the time this discussion is submitted) paper here.  While modeling greyscale images, the underlying objects are handwritten digits, an explicitly structured non-Euclidean space that we try to estimate via thet variational convolutional autoencoder, and demonstrate near-lossless, approximately Gaussian embeddings for downstream modeling and a statistical generative model that generates realistic digits -- i.e. captures the structure of the underlying manifold.
    \end{itemize}
\end{itemize}